%%%%%%%%%%%%%%%%%%%%%%%%%%%%%%%%%%%%%%%%%%%%%%%%%%%%%%%%%%%%%%%%%%%%%%%%%%%%%%%%
%%%% Resumen

\chapter*{Resumen}
%\addcontentsline{toc}{chapter}{Resumen} % si queremos que aparezca en el índice
\markboth{RESUMEN}{RESUMEN} % encabezado

WebBuilder es una plataforma web que genera automáticamente sitios web Django funcionales y desplegables a partir de los datos de cualquier API pública. El usuario introduce únicamente la URL de la API y una descripción en lenguaje natural del sitio que desea obtener (prompt), y el sistema se encarga del resto: analiza la estructura de los datos, genera el código completo del proyecto y lo despliega en un contenedor Docker accesible mediante una URL de previsualización.

El sistema integra tres tecnologías principales de forma coordinada. Django actúa como núcleo del backend, orquestando la ingestión de datos, la comunicación con los modelos de lenguaje y la generación del código. Los modelos de lenguaje (LLMs), accesibles mediante el estándar OpenAI compatible a través de diferentes proveedores, son los responsables de analizar los datos y producir el código Django adaptado a cada caso. Por último, n8n gestiona de forma asíncrona las tareas de mayor duración, como el despliegue de los contenedores Docker, el mantenimiento de los mismos, y el envío de notificaciones al usuario.

El desarrollo se ha llevado a cabo a lo largo de varios meses, partiendo de un sistema de análisis determinista y evolucionando hacia una arquitectura basada en LLM.

%%%%%%%%%%%%%%%%%%%%%%%%%%%%%%%%%%%%%%%%%%%%%%%%%%%%%%%%%%%%%%%%%%%%%%%%%%%%%%%%
%%%% Resumen en inglés

\chapter*{Summary}
%\addcontentsline{toc}{chapter}{Summary} % si queremos que aparezca en el índice
\markboth{SUMMARY}{SUMMARY} % encabezado

WebBuilder is a web platform that automatically generates functional and deployable Django websites from the data of any public API. The user only needs to provide the API URL and a natural language description of the desired site (prompt), and the system handles the rest: it analyses the data structure, generates the complete project code and deploys it in a Docker container accessible via a preview URL.
The system integrates three main technologies in a coordinated way. Django acts as the backend core, orchestrating data ingestion, communication with language models and code generation. Large language models (LLMs), accessible through the OpenAI compatible standard via different providers, are responsible for analysing the data and producing Django code tailored to each case. Finally, n8n asynchronously manages long-running tasks such as deploying and maintaining Docker containers, as well as sending notifications to the user.
Development took place over several months, starting from a deterministic analysis system and evolving into an LLM-based architecture.